\section{Conclusion}
\label{sec:conclusion}

We began with a puzzle: training data deduplication makes MMLU scores significantly
\textit{worse} ($p = 0.0114$, Bonferroni-corrected). For a practice universally prescribed
as a quality improvement, this regression demands explanation. Our work provides one:
deduplication is not a uniform quality upgrade --- it is a contamination corrector, and
MMLU was contaminated.

Our main finding is that per-benchmark accuracy differentials between dedup-Pile and
Pile models correlate positively with estimated 13-gram contamination (Pearson $r = 0.632$,
$p = 0.0086$; Spearman $\rho = 0.618$, $p = 0.0107$; $n = 16$). MMLU's Bonferroni-significant
drop ($p = 0.0114$) confirms the contamination-correction interpretation. Token-count
matching --- not step-matching --- is necessary to recover the full signal ($\Delta r = +0.093$).

Future work should test the memorization scale threshold at Pythia-6.9B, re-run the
contamination-accuracy correlation with freshly computed contamination estimates,
extend to 8--12 benchmarks to tighten statistical power, and apply the framework to
OLMo/Dolma for cross-family validation.

Training data deduplication is not a universal performance booster --- it is a
contamination corrector. Which benchmarks improve and which decline is written in
the overlap structure between the training corpus and benchmark test sets.
Reading that structure, as we demonstrate here, is now possible.
